\documentclass[11pt,a4paper]{ctexart}
\usepackage{amsmath,amssymb,mathtools,bm}
\usepackage{geometry}
\usepackage{enumitem}
\usepackage{hyperref}
\usepackage{xcolor}
\usepackage{fancyhdr}
\usepackage{titlesec}
\usepackage{booktabs,longtable,array}

\geometry{left=2.05cm,right=2.05cm,top=2.0cm,bottom=2.2cm}
\setmonofont{Consolas}
\hypersetup{
  colorlinks=true,
  linkcolor=blue!45!black,
  urlcolor=blue!45!black,
  citecolor=blue!45!black,
  pdftitle={数值代数期末备考讲义},
  pdfauthor={Codex}
}
\setlist{nosep,leftmargin=2em}
\allowdisplaybreaks
\sloppy
\pagestyle{fancy}
\setlength{\headheight}{14pt}
\fancyhf{}
\fancyhead[L]{数值代数期末备考}
\fancyhead[R]{数值代数期末备考讲义}
\fancyfoot[C]{\thepage}
\renewcommand{\headrulewidth}{0.4pt}
\newcommand{\code}[1]{\begingroup\ttfamily\small #1\endgroup}
\titlespacing*{\section}{0pt}{1.4em}{0.5em}
\titlespacing*{\subsection}{0pt}{1.0em}{0.35em}
\titlespacing*{\subsubsection}{0pt}{0.8em}{0.25em}

\title{\bfseries 数值代数期末备考讲义}
\author{基于历年卷、非上机作业与课程讲义整理}
\date{2026年6月26日}

\begin{document}
\maketitle
\tableofcontents
\newpage

\phantomsection

\section*{使用方式}

\addcontentsline{toc}{section}{使用方式}

先把每章的“识别信号”和“标准套路”背熟，再刷 \code{\detokenize{02_高频题库与预测卷.md}}。这门课的纸面考试很重视：

\begin{itemize}

\item 能不能把算法步骤写规范；

\item 能不能判断收敛、正定、可分解；

\item 能不能把证明题归约到一两个核心恒等式；

\item 能不能在小矩阵上手算不出错。

\end{itemize}

\phantomsection

\section*{总框架}

\addcontentsline{toc}{section}{总框架}

\begin{quote}

\textbf{知识主线：}

\begin{itemize}

\item 线性方程组 $Ax=b$：LU/PLU、Cholesky/$LDL^T$、追赶法。

\item 最小二乘：正则方程、QR 分解、Householder 变换。

\item 迭代法：Jacobi、Gauss-Seidel、SOR、共轭梯度法。

\item 特征值：幂法、Hessenberg 化、QR 方法、Schur 分解、SVD。

\item 误差与范数：向量/矩阵范数、条件数、扰动估计。

\end{itemize}

\end{quote}

\phantomsection

\section*{一、直接法：LU、PLU、平方根法}

\addcontentsline{toc}{section}{一、直接法：LU、PLU、平方根法}

\phantomsection

\subsection*{1. LU 分解}

\addcontentsline{toc}{subsection}{1. LU 分解}

识别信号：

\begin{itemize}

\item 题目要求 $A=LU$、$PA=LU$、$PAQ=LU$

\item 给 3 阶矩阵要求“写出详细步骤”

\item 问为什么选主元

\end{itemize}

普通 Doolittle 分解：

\[
A=LU,\quad L \text{ 单位下三角},\quad U \text{ 上三角}.
\]

第 \code{\detokenize{k}} 步：

\[
l_{ik}=\frac{a^{(k)}_{ik}}{a^{(k)}_{kk}},\quad
a^{(k+1)}_{ij}=a^{(k)}_{ij}-l_{ik}a^{(k)}_{kj}.
\]

解方程套路：

\begin{enumerate}

\item 分解 $A=LU$；

\item 解 $Ly=b$；

\item 解 $Ux=y$。

\end{enumerate}

\phantomsection

\subsection*{2. 列主元 PLU}

\addcontentsline{toc}{subsection}{2. 列主元 PLU}

列主元每步选当前列从第 \code{\detokenize{k}} 行到第 \code{\detokenize{n}} 行绝对值最大的元素作主元。标准写法：

\[
PA=LU,\quad Ax=b \Rightarrow LUx=Pb.
\]

答“为什么选主元”：

\begin{itemize}

\item 防止主元过小导致乘子过大；

\item 控制舍入误差传播；

\item 避免普通 LU 因前导主子式为零而失败；

\item 典型无普通 LU 但非奇异例子：

\end{itemize}

\[
A=\begin{pmatrix}0&1\\1&0\end{pmatrix}.
\]

\phantomsection

\subsection*{3. 高频母题：[[1,4,7],[2,5,8],[3,6,10]]}

\addcontentsline{toc}{subsection}{3. 高频母题：[[1,4,7],[2,5,8],[3,6,10]]}

\[
A=\begin{pmatrix}1&4&7\\2&5&8\\3&6&10\end{pmatrix},\quad b=(1,1,1)^T.
\]

普通 LU：

\[
L=\begin{pmatrix}1&0&0\\2&1&0\\3&2&1\end{pmatrix},\quad
U=\begin{pmatrix}1&4&7\\0&-3&-6\\0&0&1\end{pmatrix}.
\]

解：

\[
x=(-1/3,\ 1/3,\ 0)^T.
\]

一种列主元分解：

\[
P=\begin{pmatrix}0&0&1\\1&0&0\\0&1&0\end{pmatrix},
\quad
L=\begin{pmatrix}1&0&0\\1/3&1&0\\2/3&1/2&1\end{pmatrix},
\quad
U=\begin{pmatrix}3&6&10\\0&2&11/3\\0&0&-1/2\end{pmatrix}.
\]

\phantomsection

\subsection*{4. 平方根法与改进平方根法}

\addcontentsline{toc}{subsection}{4. 平方根法与改进平方根法}

适用条件：

\[
A=A^T>0.
\]

平方根法：

\[
A=LL^T.
\]

改进平方根法：

\[
A=LDL^T,
\]

其中 \code{\detokenize{L}} 为单位下三角，\code{\detokenize{D}} 为对角阵。优势是避免开方。

SPD 判别常用方法：

\begin{itemize}

\item 顺序主子式全正；

\item 对称矩阵所有特征值正；

\item 能成功 Cholesky 且对角元正。

\end{itemize}

答题模板：

\begin{enumerate}

\item 先验证 \code{\detokenize{A}} 对称；

\item 用顺序主子式或配方法证明正定；

\item 写 Cholesky/LDL\textasciicircum{}T 递推；

\item 解 $Ly=b$、$L^Tx=y$ 或 $Ly=b$、$Dz=y$、$L^Tx=z$。

\end{enumerate}

\phantomsection

\subsection*{5. 追赶法}

\addcontentsline{toc}{subsection}{5. 追赶法}

三对角矩阵：

\[
a_i x_{i-1}+b_i x_i+c_i x_{i+1}=d_i.
\]

追赶法本质是三对角 LU：

\[
l_i=\frac{a_i}{u_{i-1}},\quad u_i=b_i-l_i c_{i-1}.
\]

再前代、回代。

考试中若给三对角小矩阵，直接做普通消元也可以，但写成追赶法更对题。

\phantomsection

\section*{二、范数、条件数、误差}

\addcontentsline{toc}{section}{二、范数、条件数、误差}

\phantomsection

\subsection*{1. 向量范数}

\addcontentsline{toc}{subsection}{1. 向量范数}

必须会证明三条：

\begin{enumerate}

\item 非负性；

\item 齐次性；

\item 三角不等式。

\end{enumerate}

\code{\detokenize{p}} 范数：

\[
\|x\|_p=\left(\sum_i |x_i|^p\right)^{1/p},\quad p\ge 1.
\]

\[
\|x\|_\infty=\max_i |x_i|.
\]

常用等价关系：

\[
\|x\|_\infty \le \|x\|_2 \le \sqrt n\|x\|_\infty,
\]

\[
\|x\|_2 \le \|x\|_1 \le \sqrt n\|x\|_2,
\]

\[
\|x\|_\infty \le \|x\|_1 \le n\|x\|_\infty.
\]

\phantomsection

\subsection*{2. 矩阵范数}

\addcontentsline{toc}{subsection}{2. 矩阵范数}

矩阵范数证明多一条：

\[
\|AB\|\le \|A\|\|B\|.
\]

Frobenius 范数：

\[
\|A\|_F=\left(\sum_{i,j}|a_{ij}|^2\right)^{1/2}
=\sqrt{\operatorname{tr}(A^TA)}.
\]

证明正交不变性：

\[
\|U^TAV\|_F^2
=\operatorname{tr}(V^T A^T U U^T A V)
=\operatorname{tr}(V^T A^T A V)
=\operatorname{tr}(A^T A).
\]

相似诱导范数题：

\[
\|A\|_X=\|X^{-1}AX\|_2.
\]

证明套路：

\begin{itemize}

\item 非负、齐次、三角不等式来自 $\|.\|_2$；

\item 乘法性：

\end{itemize}

\[
\|AB\|_X=\|X^{-1}ABX\|_2
=\|(X^{-1}AX)(X^{-1}BX)\|_2
\le \|A\|_X\|B\|_X.
\]

注意：若题目严格要求矩阵范数的正定性，需要说明 \code{\detokenize{X}} 可逆，且 $X^{-1}AX=0$ 推出 $A=0$。

\phantomsection

\subsection*{3. 谱范数}

\addcontentsline{toc}{subsection}{3. 谱范数}

\[
\|A\|_2=\sqrt{\lambda_{\max}(A^TA)}.
\]

常用性质：

\[
\|A\|_2=\max_{\|x\|_2=1}\|Ax\|_2
=\max_{\|x\|_2=\|y\|_2=1}|y^TAx|.
\]

若 \code{\detokenize{U,V}} 正交：

\[
\|UA\|_2=\|AV\|_2=\|A\|_2.
\]

\phantomsection

\subsection*{4. 条件数与误差}

\addcontentsline{toc}{subsection}{4. 条件数与误差}

\[
\kappa(A)=\|A\|\|A^{-1}\|.
\]

线性方程组扰动基本估计：

\[
\frac{\|\delta x\|}{\|x\|}
\lesssim
\kappa(A)\frac{\|\delta b\|}{\|b\|}.
\]

考试解释题常用句：

残量小不等于解准确；若 \code{\detokenize{A}} 病态，前向误差可能被条件数放大。

\phantomsection

\section*{三、最小二乘与正交化}

\addcontentsline{toc}{section}{三、最小二乘与正交化}

\phantomsection

\subsection*{1. 最小二乘标准问题}

\addcontentsline{toc}{subsection}{1. 最小二乘标准问题}

\[
\min_x \|Ax-b\|_2,\quad A\in \mathbb R^{m\times n},\ m\ge n.
\]

若 \code{\detokenize{A}} 满列秩，正则方程：

\[
A^TAx=A^Tb.
\]

证明 $N(A^TA)=N(A)$：

\begin{itemize}

\item 若 $Ax=0$，显然 $A^TAx=0$；

\item 若 $A^TAx=0$，左乘 $x^T$ 得

\end{itemize}

\[
0=x^TA^TAx=\|Ax\|_2^2,
\]

所以 $Ax=0$。

\phantomsection

\subsection*{2. 投影矩阵}

\addcontentsline{toc}{subsection}{2. 投影矩阵}

满列秩时：

\[
P=A(A^TA)^{-1}A^T.
\]

性质：

\[
P^T=P,\quad P^2=P,\quad \|P\|_2=1
\]

只要 \code{\detokenize{P}} 不是零矩阵。证明二范数为 1 的套路：

\begin{itemize}

\item \code{\detokenize{P}} 对称幂等，所以特征值只可能是 0 或 1；

\item \code{\detokenize{A}} 满列秩且 $n>0$，\code{\detokenize{P}} 有特征值 1；

\item 因此 $\|P\|_2=max |lambda_i|=1$。

\end{itemize}

\phantomsection

\subsection*{3. QR 解最小二乘}

\addcontentsline{toc}{subsection}{3. QR 解最小二乘}

若

\[
A=Q\begin{pmatrix}R\\0\end{pmatrix},\quad Q=[Q_1,Q_2],
\]

则

\[
\|Ax-b\|_2^2
=\|Rx-Q_1^Tb\|_2^2+\|Q_2^Tb\|_2^2.
\]

所以解

\[
Rx=Q_1^Tb.
\]

\phantomsection

\subsection*{4. Householder}

\addcontentsline{toc}{subsection}{4. Householder}

定义：

\[
H=I-2ww^T,\quad \|w\|_2=1.
\]

性质：

\[
H^T=H,\quad H^TH=I,\quad H^2=I.
\]

构造 $Hx=\alpha e_1$：

\[
\alpha=\pm \|x\|_2,\quad
w=\frac{x-\alpha e_1}{\|x-\alpha e_1\|_2}.
\]

实算建议：选符号避免相消。

高频母题：

\[
x=(3,1,6,4,2,2)^T,\quad Hx=(3,\alpha,4,6,0,0)^T.
\]

只需作用在第 2、5、6 个分量组成的向量

\[
y=(1,2,2)^T.
\]

取 $\alpha=3$，可用

\[
H_3=
\begin{pmatrix}
1/3&2/3&2/3\\
2/3&1/3&-2/3\\
2/3&-2/3&1/3
\end{pmatrix},
\quad H_3y=(3,0,0)^T.
\]

再把 $H_3$ 嵌入到第 2、5、6 坐标上即可。

\phantomsection

\subsection*{5. Givens}

\addcontentsline{toc}{subsection}{5. Givens}

二维旋转：

\[
G=
\begin{pmatrix}
c&s\\
-s&c
\end{pmatrix},\quad c^2+s^2=1.
\]

若要把 $(a,b)^T$ 化为 $(r,0)^T$，取

\[
c=\frac a{\sqrt{a^2+b^2}},\quad s=\frac b{\sqrt{a^2+b^2}}.
\]

若题目要求变为 $\beta(1,1)^T$，列方程：

\[
ca+sb=\beta,\quad -sa+cb=\beta,\quad c^2+s^2=1.
\]

高频母题：

\[
\begin{pmatrix}c&s\\-s&c\end{pmatrix}
\begin{pmatrix}7\\-1\end{pmatrix}
=\begin{pmatrix}\beta\\\beta\end{pmatrix}.
\]

可取

\[
c=3/5,\quad s=-4/5,\quad \beta=5.
\]

\phantomsection

\section*{四、定常迭代法}

\addcontentsline{toc}{section}{四、定常迭代法}

\phantomsection

\subsection*{1. 统一分裂}

\addcontentsline{toc}{subsection}{1. 统一分裂}

常用约定：

\[
A=D-L-U.
\]

Jacobi：

\[
x^{(k+1)}=D^{-1}(L+U)x^{(k)}+D^{-1}b.
\]

G-S：

\[
x^{(k+1)}=(D-L)^{-1}Ux^{(k)}+(D-L)^{-1}b.
\]

SOR：

\[
x^{(k+1)}
=(D-\omega L)^{-1}[(1-\omega)D+\omega U]x^{(k)}
+\omega(D-\omega L)^{-1}b.
\]

收敛充要条件：

\[
\rho(B)<1.
\]

\phantomsection

\subsection*{2. 常用充分条件}

\addcontentsline{toc}{subsection}{2. 常用充分条件}

严格对角占优：

\[
|a_{ii}|>\sum_{j\ne i}|a_{ij}|
\]

可推出 Jacobi、G-S 收敛。

若 \code{\detokenize{A}} 对称正定，G-S 收敛；SOR 对 SPD 且 $0<\omega<2$ 收敛。

Homework10 还强调：

\begin{itemize}

\item 严格或弱严格对角占优不可约矩阵，G-S 收敛；

\item 若松弛因子 $\omega in (0,1]$，在相应条件下松弛迭代收敛；

\item \code{\detokenize{2x2}} SPD 可推出 Jacobi 收敛。

\end{itemize}

\phantomsection

\subsection*{3. 参数矩阵模板}

\addcontentsline{toc}{subsection}{3. 参数矩阵模板}

\[
A=\begin{pmatrix}1&0&a\\0&1&0\\a&0&1\end{pmatrix}.
\]

正定：

\[
1-a^2>0 \Rightarrow |a|<1.
\]

Jacobi 迭代矩阵：

\[
B_J=
\begin{pmatrix}0&0&-a\\0&0&0\\-a&0&0\end{pmatrix},
\]

特征值为 \code{\detokenize{0, a, -a}}，所以收敛当且仅当 $|a|<1$。

G-S 迭代矩阵特征值为 $0,0,a^2$，所以同样 $|a|<1$。

\phantomsection

\subsection*{4. 3 阶矩阵写 Jacobi 迭代矩阵}

\addcontentsline{toc}{subsection}{4. 3 阶矩阵写 Jacobi 迭代矩阵}

给 \code{\detokenize{A}} 后直接写：

\[
B_J=-D^{-1}(L_A+U_A),
\]

其中 $L_A+U_A=A-D$ 是原矩阵非对角部分。

2024-2025 题中：

\[
A_1=\begin{pmatrix}2&1&1\\1&-1&-1\\1&1&-2\end{pmatrix},
\]

\[
B_J=
\begin{pmatrix}
0&-1/2&-1/2\\
1&0&-1\\
1/2&1/2&0
\end{pmatrix},
\]

特征值为 \code{\detokenize{0, ± i sqrt(5)/2}}，谱半径 $sqrt(5)/2>1$，不收敛。

\[
A_2=\begin{pmatrix}1&2&-2\\1&1&-1\\2&-2&1\end{pmatrix},
\]

\[
B_J=
\begin{pmatrix}
0&-2&2\\
-1&0&1\\
-2&2&0
\end{pmatrix},
\]

特征值全为 0，收敛。

\phantomsection

\section*{五、共轭梯度法}

\addcontentsline{toc}{section}{五、共轭梯度法}

适用条件：

\[
A=A^T>0.
\]

基本概念：

\begin{itemize}

\item $p_i^T A p_j=0$ 称为 A-共轭；

\item 残量 $r_k=b-Ax_k$；

\item 在精确算术下，CG 至多 \code{\detokenize{n}} 步得到精确解。

\end{itemize}

\phantomsection

\subsection*{1. A-共轭向量线性无关}

\addcontentsline{toc}{subsection}{1. A-共轭向量线性无关}

证明套路：

设

\[
\sum_i c_i p_i=0.
\]

左乘 $p_j^T A$：

\[
c_j p_j^T A p_j=0.
\]

由于 \code{\detokenize{A}} SPD，$p_j^T A p_j>0$，故 $c_j=0$。

\phantomsection

\subsection*{2. 逆矩阵展开式}

\addcontentsline{toc}{subsection}{2. 逆矩阵展开式}

若 $p_1,...,p_n$ 是 A-共轭基，则

\[
A^{-1}=\sum_{k=1}^n
\frac{p_kp_k^T}{p_k^T A p_k}.
\]

证明套路：

任意 \code{\detokenize{x}} 可按 $p_k$ 展开：

\[
x=\sum_k \alpha_k p_k.
\]

由 A-正交性得

\[
\alpha_k=\frac{p_k^T A x}{p_k^T A p_k}.
\]

令 $x=A^{-1}y$，得到

\[
A^{-1}y=\sum_k
\frac{p_k p_k^T y}{p_k^T A p_k}.
\]

任意 \code{\detokenize{y}} 成立，即矩阵等式成立。

\phantomsection

\subsection*{3. 小矩阵 CG 手算}

\addcontentsline{toc}{subsection}{3. 小矩阵 CG 手算}

考试若要求“用共轭梯度法求解”，要写：

\[
\alpha_k=\frac{r_k^Tr_k}{p_k^TAp_k},\quad
x_{k+1}=x_k+\alpha_k p_k,
\]

\[
r_{k+1}=r_k-\alpha_k Ap_k,\quad
\beta_k=\frac{r_{k+1}^Tr_{k+1}}{r_k^Tr_k},\quad
p_{k+1}=r_{k+1}+\beta_kp_k.
\]

常见证明题：

\[
(r^{(k)},r^{(k)})=-\alpha_{k-1}(r^{(k)},Ap^{(k-1)}),
\]

\[
(r^{(k)},r^{(k)})=\alpha_k(p^{(k)},Ap^{(k)}).
\]

思路是把 $r_{k+1}=r_k-alpha_k Ap_k$ 和 $p_k=r_k+beta_{k-1}p_{k-1}$ 代入，再用残量正交、方向共轭。

\phantomsection

\section*{六、特征值问题}

\addcontentsline{toc}{section}{六、特征值问题}

\phantomsection

\subsection*{1. 幂法}

\addcontentsline{toc}{subsection}{1. 幂法}

若

\[
|\lambda_1|>|\lambda_2|\ge \cdots
\]

且初始向量在主特征向量方向分量非零，则幂法收敛到主特征向量方向。

异常模板：

\[
A=\begin{pmatrix}\lambda&1\\0&\lambda\end{pmatrix}
\]

只有一个特征向量，幂法会带 Jordan 多项式因子，收敛速度和普通对角化情形不同。

\[
B=\begin{pmatrix}\lambda&1\\0&-\lambda\end{pmatrix}
\]

有两个等模特征值，迭代可能振荡，不满足唯一模最大特征值条件。

\phantomsection

\subsection*{2. Householder 上 Hessenberg 化}

\addcontentsline{toc}{subsection}{2. Householder 上 Hessenberg 化}

目标：

\[
Q^TAQ=H,
\]

其中 \code{\detokenize{H}} 为上 Hessenberg，即 $h_{ij}=0$ 当 $i>j+1$。

算法描述：

第 \code{\detokenize{k}} 步取第 \code{\detokenize{k}} 列的子向量

\[
x=A_{k+1:n,k},
\]

构造 Householder $P_k$ 使其尾部只剩第一个分量。嵌入为

\[
Q_k=\operatorname{diag}(I_k,P_k),
\]

做相似变换：

\[
A\leftarrow Q_k^T A Q_k.
\]

\phantomsection

\subsection*{3. QR 迭代}

\addcontentsline{toc}{subsection}{3. QR 迭代}

基本迭代：

\[
A_k=Q_kR_k,\quad A_{k+1}=R_kQ_k=Q_k^TA_kQ_k.
\]

带位移：

\[
A_k-\mu_k I=Q_kR_k,\quad
A_{k+1}=R_kQ_k+\mu_k I.
\]

证明保持 Hessenberg：

\begin{itemize}

\item Hessenberg 矩阵的 QR 分解可由相邻 Givens 消去完成；

\item 这些 Givens 只在相邻行列作用；

\item 左乘消去次对角以下元素，右乘不会在 Hessenberg 带宽外产生新非零元；

\item 因而每一步仍是上 Hessenberg。

\end{itemize}

\phantomsection

\subsection*{4. Schur 分解}

\addcontentsline{toc}{subsection}{4. Schur 分解}

实/复版本要分清：

复 Schur：

\[
A=QTQ^*,\quad Q \text{ 酉},\quad T \text{ 上三角}.
\]

实 Schur：

\[
A=QTQ^T,
\]

\code{\detokenize{T}} 为拟上三角，对角块为 \code{\detokenize{1x1}} 或 \code{\detokenize{2x2}}。

考试若要求“默写实数域上的 Schur 分解定理”，写实 Schur。

\phantomsection

\subsection*{5. SVD 与 Moore-Penrose 逆}

\addcontentsline{toc}{subsection}{5. SVD 与 Moore-Penrose 逆}

SVD：

\[
A=U\Sigma V^T,\quad
\sigma_1\ge \cdots\ge \sigma_r>0.
\]

Moore-Penrose 逆：

\[
A^+=V\Sigma^+U^T.
\]

常见不等式：

\[
\sigma_1\|v\|_2\ge \|Av\|_2\ge \sigma_n\|v\|_2
\]

满列秩时投影：

\[
AA^+=A(A^TA)^{-1}A^T.
\]

若题目说 $x=Xb$ 对任意 \code{\detokenize{b}} 都最小化 $\|b-Ax\|_2$，则 \code{\detokenize{AX}} 是到 \code{\detokenize{R(A)}} 的正交投影，所以：

\[
AXA=A,\quad (AX)^T=AX.
\]

\phantomsection

\section*{七、证明题统一套路}

\addcontentsline{toc}{section}{七、证明题统一套路}

\phantomsection

\subsection*{范数证明}

\addcontentsline{toc}{subsection}{范数证明}

按四条写：

\begin{enumerate}

\item 非负性和零向量；

\item 齐次性；

\item 三角不等式；

\item 矩阵范数再加次乘性。

\end{enumerate}

\phantomsection

\subsection*{收敛证明}

\addcontentsline{toc}{subsection}{收敛证明}

三种入口：

\begin{enumerate}

\item 直接算迭代矩阵 \code{\detokenize{B}}，看 $\rho(B)<1$；

\item 用严格对角占优或 SPD 充分条件；

\item 参数题求特征值，解不等式。

\end{enumerate}

\phantomsection

\subsection*{正定证明}

\addcontentsline{toc}{subsection}{正定证明}

三种入口：

\begin{enumerate}

\item 顺序主子式；

\item 配方 $x^TAx>0$；

\item 特征值全正。

\end{enumerate}

\phantomsection

\subsection*{最小二乘证明}

\addcontentsline{toc}{subsection}{最小二乘证明}

三种入口：

\begin{enumerate}

\item 正交投影：残差垂直于 \code{\detokenize{R(A)}}；

\item 正则方程：$A^T(Ax-b)=0$；

\item QR 分解：把残差分成可控项和不可控项。

\end{enumerate}

\phantomsection

\subsection*{正交变换证明}

\addcontentsline{toc}{subsection}{正交变换证明}

牢记：

\[
Q^TQ=I \Rightarrow \|Qx\|_2=\|x\|_2.
\]

Householder/Givens 都是围绕这个式子展开。

\phantomsection

\section*{八、常见失分点}

\addcontentsline{toc}{section}{八、常见失分点}

\begin{itemize}

\item $PA=LU$ 和 $A=PLU$ 的记号混用。卷面上先声明自己的约定。

\item PLU 解方程时漏掉 \code{\detokenize{Pb}}。

\item 证明矩阵范数时漏掉 $\|AB\|<=\|A\|\|B\|$。

\item 最小二乘正则方程误写成 $AA^Tx=Ab$。

\item Cholesky 用在非对称或非正定矩阵上。

\item Jacobi/G-S 收敛性只看对角占优，不会算谱半径；参数题最好直接求特征值。

\item Householder 的 \code{\detokenize{H}} 左右相似变换用于特征值问题，单侧左乘用于 QR/最小二乘。

\item QR 迭代题要写“相似变换”，否则说不清特征值为什么不变。

\item 残量小不等于误差小，病态矩阵必须提条件数。

\end{itemize}

\end{document}
